\scanpage{16}
Soient $V$ un Module \deleted{loc.} libre \inserted{de rang $2$} sur un \deleted{topos loc. annelé}\inserted{anneau $\Lambda$}, d'où $P(V)=$ ens des ``droites vectorielles'' dans $V$ (i.e des sous modules \inserted{\deleted{\unsure{loc.}}} \deleted{libres \unsure{d|de}} inversibles facteurs directs),

\deleted{et $P(V)^{\unsure{'|!}}=$}

\[
\begin{aligned}
V^*&=\text{ens. des éléments de }V\text{ partout }\ne0\quad(\Longleftrightarrow\text{ qui}\\
&\hspace{2em}\text{font partie d'une base})
\end{aligned}
\tag*{(1)}
\]
\[
\begin{aligned}
P(V)^!&=V^*/\Lambda^*\hookrightarrow P(V)\quad\text{l'ens des ``droites dans }V\text{''}\\
&\text{qui sont des modules libres de rang }1.
\end{aligned}
\tag*{(2)}
\]
On a envie d'expliciter comment on récupère \inserted{le $\Lambda$-module} $V$ à l'aide de l'ens $P(V)^!$, muni d'une structure supplémentaire, qui est \deleted{un} essentiellement le $\Lambda^*$-\deleted{\unsure{fibré}}\unsure{torseur} relatif (ensembliste) $V^*$ sur $P(V)^!=X$. Pour $x\in P(V)^!$, soit $D_x\subset V$ la droite vectorielle correspondante, et posons
\[
T_x=D_{\unsure{a|x}}^*\quad(\deleted{\unsure{A}}\Lambda^*\text{-torseur})
\tag*{(3)}
\]
Soit de plus
\[
\omega_V=\det(V)^*\qquad(\text{un }\Lambda^{\unsure{*|\alpha|\circ}}\text{-torseur}
\tag*{(4)}
\]
Sur $P(V)^!$ on a la relation binaire de disjonction
\[
x,y\in P(V)^!\quad\text{``\uline{disjoints}''}\quad\text{\unsure{ass|ssi}}\quad V\simeq D_x\oplus D_y
\]
(si $\Lambda$ est un corps, cela équivaut à $x\ne y$; si $\Lambda$ est local, cela équivaut à $x_0\ne y_0$ (l'indice $0$
% archive sheet number: 15
